\section{Related work and the AGB-orthogonality stance}
\label{sec:related}

The neighbourhood DynamicMCPBench sits in:
\textbf{AgentGraphBench (AGB)} for graph-shaped tool-use evaluation;
\textbf{MCPEval} for generate-then-verify task synthesis over MCP
servers; \textbf{MCP-Bench} and \textbf{MCP-Atlas} for substrate
breadth; \textbf{$\tau$-bench} and \textbf{ToolSandbox} for
state-coupled tool use; \textbf{StableToolBench} for caching as a
determinism mechanism.

This paper's structural choice is to \textbf{stay orthogonal to AGB}
along the four pillars below. The headline AGB connective thesis we
adopt is that they \emph{diagnosed the disease} (GT tool lists are
$\sim$50\,\% noise) and that DynamicMCPBench offers a \emph{different
cure} (trace-grounding makes the unnecessary-tool problem impossible by
construction~--- every tool in a trace was actually invoked).

\begin{table*}[t]
  \centering
  \small
  \begin{tabular}{lll}
    \toprule
    dimension & AgentGraphBench (AGB) & DynamicMCPBench (this paper) \\
    \midrule
    Organizing object & tool-dependency graph + motifs & \textbf{execution trace} (real trajectory) \\
    Generation direction & backward: subgraph $\to$ back-instruct & \textbf{forward: explore live $\to$ distill} \\
    Substrate & static cached catalog (StableToolBench-style) & \textbf{live, stateful, crawled MCP servers} \\
    Ground truth & GT tool set (denoised) & \textbf{reference trace $\to$ effect checkpoints} \\
    Eval signal & tool-selection recall on a candidate list & \textbf{trace/effect alignment + state} \\
    Final answer & checked & \textbf{deliberately not checked} \\
    Difficulty axis & graph-motif class + subgraph size & \textbf{trace complexity + dynamism level} \\
    \bottomrule
  \end{tabular}
  \caption{Four pillars along which DynamicMCPBench stays orthogonal to AgentGraphBench.}
  \label{tab:agb_orthogonality}
\end{table*}

We pre-empt three reviewer objections from the AGB neighbourhood:

\begin{enumerate}
  \item \emph{``Isn't this just AGB without the graph?''} No~--- removing
    the graph forces a different \emph{generator} (forward exploration),
    a different \emph{ground truth} (trace, not tool set), and a
    different \emph{evaluation} (effect alignment, not recall). The
    graph re-enters this paper \emph{only} as one of two clearly-labeled
    RQ2 baselines (\S\ref{sec:rq2}).
  \item \emph{``Isn't this just MCPEval?''} No~--- MCPEval is
    generate-then-verify; we are \emph{explore-then-distill}. The trace
    is the primitive, not a verification afterthought; we score effects
    and state, not tool-call matching against a generated plan; we
    refuse to grade the final answer.
  \item \emph{``Traces have multiple valid paths~--- is that fair?''}
    Yes, by design. The reference trace compiles into checkpoints with
    equivalence sets, causal ordering \emph{only where a real
    dependency exists}, and minefields. Any trajectory satisfying the
    checkpoints without tripping a minefield passes
    (\S\ref{sec:method:eval}).
\end{enumerate}
